\chapter*{Résumé}
\addcontentsline{toc}{chapter}{Résumé}

\textbf{Mots-clés}~: homologation ANSSI, SOC, Wazuh, Active
Directory, audit drift, evidence-as-code, automatisation,
pipeline Python.

\bigskip

Une homologation ANSSI s'appuie sur un dossier daté et signé. Le
problème, c'est que le SI continue à bouger après la signature~:
règles réseau modifiées, nouvelles CVE, comptes Active Directory
qui changent, règles SIEM qui dérivent. Au bout de quelques mois,
le dossier ne décrit plus ce qui tourne réellement.

Ce mémoire présente deux choses concrètes~:

\begin{itemize}
\item \textbf{Un lab d'expérimentation} appelé UrbanUpC, monté sur
  Azure~: trois machines virtuelles (un contrôleur de domaine
  Windows Server, un serveur web Linux exposant deux applications
  PHP volontairement vulnérables, un SIEM Wazuh), un Active
  Directory \texttt{corpnet.local} avec utilisateurs et groupes
  réalistes. Le lab a généré 927 alertes Wazuh sur des scénarios
  d'attaque web et AD documentés.

\item \textbf{Un outil appelé HOMO-CI}, qui collecte automatiquement
  les preuves techniques nécessaires à l'homologation (règles
  NSG, CVE des images Docker, inventaire AD, règles Wazuh) et
  regénère le dossier d'homologation au format ANSSI à chaque
  exécution. C'est un pipeline Python avec des watchers
  spécialisés, un stockage SQLite, un moteur de diff, et un
  rendu Jinja2 vers Markdown puis PDF signé.
\end{itemize}

L'objectif n'est pas de remplacer l'homologation officielle, mais
de \emph{réduire l'écart} entre le dossier et la réalité --- ce
qu'on appelle dans la suite l'\emph{audit drift}.

Les tests sur le lab donnent trois résultats~: l'outil couvre une
partie significative des preuves techniques attendues, propage un
changement Azure en quelques secondes, et tient un dossier à jour
là où la version statique se périme. Les chiffres précis et les
limites honnêtes sont en partie~IV.

\chapter*{Abstract}
\addcontentsline{toc}{chapter}{Abstract}

\textbf{Keywords}~: ANSSI accreditation, SOC, Wazuh, Active
Directory, audit drift, evidence-as-code, automation, Python
pipeline.

\bigskip

An ANSSI accreditation relies on a dated, signed dossier. The
problem is that the IS keeps moving after signature: firewall
rules change, new CVEs appear, Active Directory accounts shift,
SIEM rules drift. After a few months, the dossier no longer
describes what is actually running.

This thesis presents two concrete things:

\begin{itemize}
\item \textbf{An experimentation lab} called UrbanUpC, built on
  Azure: three virtual machines (a Windows Server domain
  controller, a Linux web server hosting two intentionally
  vulnerable PHP apps, a Wazuh SIEM), an Active Directory
  \texttt{corpnet.local} with realistic users and groups. The lab
  produced 927 Wazuh alerts on documented web and AD attack
  scenarios.

\item \textbf{A tool called HOMO-CI}, which automatically collects
  the technical evidence needed for accreditation (NSG rules,
  Docker image CVEs, AD inventory, Wazuh rules) and regenerates
  the ANSSI-format dossier on each run. It is a Python pipeline
  with specialised watchers, SQLite storage, a diff engine, and
  Jinja2 rendering to Markdown then signed PDF.
\end{itemize}

The goal is not to replace the official accreditation, but to
\emph{reduce the gap} between the dossier and reality --- what we
call \emph{audit drift}.

Lab tests give three results: the tool covers a meaningful share
of the expected technical evidence, propagates an Azure change in
seconds, and keeps the dossier fresh where the static version
goes stale. Exact numbers and honest limits are in part~IV.
